% Automorphisms and equivalence classes of the order-18 Jacobi Hadamard matrices
% Compile with pdflatex (see the repository Makefile).

\documentclass[11pt]{article}
\usepackage[T1]{fontenc}
\usepackage[margin=1in]{geometry}
\usepackage{amsmath,amssymb,amsthm}
\usepackage{booktabs}
\usepackage[colorlinks=true,linkcolor=blue,citecolor=blue,urlcolor=blue]{hyperref}

\theoremstyle{plain}
\newtheorem{theorem}{Theorem}[section]
\newtheorem{proposition}[theorem]{Proposition}
\newtheorem{corollary}[theorem]{Corollary}
\newtheorem{lemma}[theorem]{Lemma}
\theoremstyle{definition}
\newtheorem{definition}[theorem]{Definition}
\theoremstyle{remark}
\newtheorem{remark}[theorem]{Remark}

\newcommand{\F}{\mathbb F}
\newcommand{\Z}{\mathbb Z}
\newcommand{\Q}{\mathbb Q}
\newcommand{\abs}[1]{\lvert#1\rvert}
\newcommand{\Aut}{\mathrm{Aut}}
\newcommand{\Gal}{\mathrm{Gal}}
\newcommand{\sgn}{\mathrm{sgn}}
\setlength{\emergencystretch}{3em}
\hypersetup{
  pdftitle={Automorphisms and equivalence classes of the order-18 Jacobi Hadamard matrices},
  pdfauthor={Jeffery Kline}
}

\title{Automorphisms and equivalence classes of the\\ order-18 Jacobi Hadamard matrices}
\author{Jeffery Kline}
\date{November 17, 2026\\[3pt]\small Version 0.1.0, \href{https://doi.org/10.5281/zenodo.24329143}{doi:10.5281/zenodo.24329143}}

\begin{document}
\maketitle

\begin{abstract}
An earlier paper constructed skew Hadamard matrices of orders
$4(343^d+1)$, for every odd $d$, by placing three sign patterns on the
eighteen classes of an order-18 multiplicative character of the field with
$343^d$ elements, together with the Paley pattern, in the bordered
Wallis--Whiteman array. That paper proved that the patterns work at every
order and left the equivalence of the completed matrices open. Here we
determine their automorphism groups. Under independent signed permutations
of rows and columns, every automorphism fixes the border and acts on each of
the eight copies of the field by the same map $x\mapsto ax^{343^j}$, where
$a$ is an eighteenth power; these maps form a group of order $d(343^d-1)/18$, which is
$19$ at the first order, $1376$. The proof recovers the border, the field
blocks and their origins from the matrix alone: exactly two sets of four rows
have a product that is constant away from four columns, the Paley
submatrices are recognized by their row correlations, and Kantor's theorem
on the automorphisms of the Paley design makes the coordinate maps
semilinear. We give a second set of three patterns and prove that both sets
work for every prime $p\equiv7$ or $31\pmod{36}$ and field size $p^{3d}$ with $d$ odd.
The same invariants show that the resulting matrices give at least eight
inequivalent skew Hadamard matrices at every such order $4(p^{3d}+1)$. On the
original tower the two sets, their arrangements in the array and the
transposes give exactly sixteen equivalence classes when $3\nmid d$, and
fourteen or sixteen when $3\mid d$. None of these matrices is equivalent to
a Kronecker product, to a skew doubling, or to a matrix in the same array
whose patterns are constant on the classes of sixth powers. All of these
orders were known to admit skew Hadamard matrices.
\end{abstract}

\section{Introduction}\label{sec:intro}

\subsection{The matrices}
A Hadamard matrix of order $N$ is an $N\times N$ matrix $H$ with entries
$\pm1$ and $HH^{\mathsf T}=NI$; it is skew if $H+H^{\mathsf T}=2I$. Two
Hadamard matrices are \emph{equivalent} if one is obtained from the other by
permuting rows, permuting columns, and multiplying rows and columns by $-1$,
all independently. An \emph{automorphism} of $H$ is such an operation that
maps $H$ to itself. Equivalence and automorphisms are the standard way to
decide whether two constructions are really different.

The earlier paper~\cite{Kline18} worked in the field $E=\F_Q$ with
$Q=343^d$, $d$ odd. It took an order-18 multiplicative character of $E$,
three fixed sign patterns on its eighteen classes, and the Paley (quadratic
character) pattern, and placed the four resulting $\pm1$ functions in the
bordered array of Wallis and Whiteman~\cite{WallisWhiteman}. It proved that
the result is a skew Hadamard matrix of order $4(Q+1)$ for every odd $d$,
beginning with order $1376$, and that the three patterns are not the
classical sextic patterns in disguise. It did not decide whether the
completed matrices are new up to equivalence.

\subsection{Results}
Write $\Lambda_Q=d(Q-1)/18$ when $Q=p^{3d}$. Section~\ref{sec:array} defines the
array precisely, and the two pattern systems I and II.

\begin{theorem}[Automorphisms]\label{thm:aut}
Let $Q=343^d$ with $d$ odd, and let $H$ be the matrix of~\cite{Kline18}.
Every automorphism of $H$ fixes the four border rows and the four border
columns, maps each of the eight copies of $E$ to itself, and acts on all of
them by one map
\[
 x\longmapsto a\,x^{343^j},\qquad a\in(E^\times)^{18},\ j\in\Z/d\Z,
\]
together with, possibly, negation of every row and every column. Every such
map is an automorphism. The group of these permutations has order
$\Lambda_Q=d(Q-1)/18$; it is cyclic of order $19$ when $Q=343$.
\end{theorem}

\begin{theorem}[Two systems in infinitely many characteristics]\label{thm:systems}
Let $p$ be a prime with $p\equiv7$ or $31\pmod{36}$, let $d$ be odd, put
$Q=p^{3d}$ and fix any character of order $18$ of $\F_Q^\times$. For each of
the systems I and II of Table~\ref{tab:words}, the bordered array gives a
skew Hadamard matrix of order $4(Q+1)$. The matrices obtained from the two
systems, with or without the first pattern reversed, and their transposes,
are eight pairwise inequivalent Hadamard matrices.
\end{theorem}

\begin{theorem}[Sixteen classes]\label{thm:sixteen}
Let $Q=343^d$ with $d$ odd and use the character of~\cite{Kline18}. Arrange
the four patterns of system I or of system II in the four positions of the
array in any order, and reverse any subset of them. These $768$ matrices and
their transposes fall into exactly $16$ equivalence classes if $3\nmid d$,
and into $14$ or $16$ classes if $3\mid d$. In particular there are exactly
$16$ at order $1376$.
\end{theorem}

\begin{theorem}[Comparisons]\label{thm:compare}
None of the matrices in Theorems~\ref{thm:systems} and~\ref{thm:sixteen} is
equivalent to its transpose, to a Kronecker product of two Hadamard matrices
of orders greater than one, or to a skew doubling
$\begin{pmatrix}K&K\\-K^{\mathsf T}&K^{\mathsf T}\end{pmatrix}$ of a skew
Hadamard matrix $K$. None is equivalent to a Hadamard matrix of the same
bordered array whose four functions are constant on the cosets of the
sixth powers.
\end{theorem}

All orders $4(p^{3d}+1)$ here were already known to have skew Hadamard
matrices, for example from the Paley matrix by doubling. The theorems
concern the matrices, not the orders.

\subsection{How the proofs work}
The array has four border rows and columns and four copies of the field on
each side. The proofs read this structure off the matrix itself, so that it
must be preserved by any equivalence.

\emph{Four-row sets} (Section~\ref{sec:quartets}). For four rows, multiply
them entrywise and count the entries in the minority sign. Exactly two sets
of four rows have at most four minority entries: the border rows and the
four rows at the origin of each field copy. For large fields this follows
from Weil's bound for multiplicative character sums; at $Q=343$ an exhaustive
search confirms it. This invariant already separates the matrices from
Kronecker products, skew doublings and transposes.

\emph{Paley blocks} (Section~\ref{sec:frame}). Signs on the border columns
then split the remaining rows into the four field copies, and dually for
columns. Among the sixteen field blocks, the four Paley blocks are the ones
in which every two rows have inner product $\pm1$. Kantor's theorem on the
automorphisms of the Paley design~\cite{Kantor2T} then forces every
coordinate map to be semiaffine, and the origins fix the translations.

\emph{Synchronization} (Section~\ref{sec:frame}). Each pattern sums to $1$
over the field, so it has no additive period; this forces the eight
coordinate maps to be one map $x\mapsto\beta x^{p^k}$. What remains is a
finite problem over the eighteen classes, which a computer enumerates.

\emph{Telling the border from the origin rows} (Sections~\ref{sec:aut}
and~\ref{sec:classes}). The argument so far applies to equivalences that
send border rows to border rows. In characteristic $7$ an exact
calculation shows that only the border marking produces Paley blocks, at
every odd $d$, by reducing cyclotomic numbers modulo $37$ and lifting them
from $\F_{343}$. For one pair of arrangements of system II a finer
histogram, computed modulo $19$, is needed; it decides the question when
$3\nmid d$.

\subsection{Earlier work and status}
The bordered array is due to Wallis and Whiteman~\cite{WallisWhiteman}.
Hall~\cite{Hall62} found a large group of automorphisms of the Paley
Hadamard matrix and the whole group at order $12$, Kantor~\cite{Kantor69}
determined the whole group for $q>11$, and Kantor~\cite{Kantor2T} determined
that of the Paley design; we use the latter. Lehmer~\cite{Lehmer74} constructed
supplementary difference sets from cyclotomy, and Yamada~\cite{Yamada92}
related such sets to Jacobi sums. Products of four rows are a classical tool
for distinguishing Hadamard matrices; Orrick~\cite{Orrick} uses closed
quadruples of rows, the closed quartets of Section~\ref{sec:quartets}, in his
switching operations. Momihara and Xiang~\cite{MomiharaXiang} constructed
skew Hadamard difference families with $2^{u-1}$ blocks in $\F_{q^e}$ for
$q\equiv2^u+1\pmod{2^{u+1}}$, $u\ge2$, and every $e$, from cyclotomy of order
a power of $2$; these congruences force $q\equiv1\pmod4$, so none of their
families lives in a field $\F_Q$ of this paper, all of which have
$Q\equiv3\pmod4$. Pantangi~\cite{Pantangi}
computed critical groups of the tournaments attached to one-, two- and
four-block families and used them to show that those skew Hadamard matrices
are not equivalent to the Paley matrix. Karoui \cite{Karoui}
gave a skew Hadamard matrix of order $1252$ from order-16 cyclotomic unions
and exhibited an explicit subgroup of its automorphism group. The
construction of system I in characteristic $7$ is from~\cite{Kline18};
system II and the extension to $p\equiv7,31\pmod{36}$ are new here, as far as
we know. We found the theorems above in none of the sources we checked; the
search is recorded in the repository and is not a claim of global priority.

\section{The array and the two systems}\label{sec:array}

\subsection{Patterns}
Let $p\equiv7$ or $31\pmod{36}$ be prime, $d$ odd and $Q=p^{3d}$. Then
$p^3\equiv19\pmod{36}$, so $Q\equiv19\pmod{36}$; in particular
$Q\equiv3\pmod4$ and $18\mid Q-1$ with $(Q-1)/18$ odd. Let $\chi$ be a
character of order $18$ of $E^\times$, $E=\F_Q$, and for $x\neq0$ let
$j(x)\in\Z/18\Z$ be defined by $\chi(x)=e^{2\pi ij(x)/18}$. Since
$(Q-1)/18$ is odd, $j(-1)=9$.

A \emph{half-word} is $w=(w_0,\dots,w_8)\in\{\pm1\}^9$; extend it by
$w_{j+9}=-w_j$. Its \emph{pattern} is the function
\[
 f_w(x)=w_{j(x)}\quad(x\neq0),\qquad f_w(0)=1.
\]
Because $j(-1)=9$, every pattern is odd away from the origin:
$f_w(-x)=-f_w(x)$ for $x\neq0$. The pattern of $w_D=(+,-,+,-,+,-,+,-,+)$ is
the quadratic character with value $1$ at $0$; we call it the Paley pattern.
Each pattern takes the values $\pm1$ equally often on $E^\times$, so
$\sum_{x\in E}f_w(x)=1$.

\begin{table}[ht]
\centering
\begin{tabular}{lccc}
\toprule
System & $w_A$ & $w_B$ & $w_C$\\
\midrule
I  & $+{-}{-}{-}{-}{-}{-}{-}+$ & $+{-}{-}+{-}{-}{-}+{-}$ & $+{-}+++{-}++{-}$\\
II & $+{-}+++++++$ & $++++++{-}+{-}$ & $+++{-}+{-}{-}{-}{-}$\\
\bottomrule
\end{tabular}
\caption{The half-words of the two systems. The fourth pattern of each
system is the Paley pattern $w_D=+{-}+{-}+{-}+{-}+$.}\label{tab:words}
\end{table}

For the characteristic-$7$ tower, $Q=343^d$, the character
of~\cite{Kline18} is $\chi=\lambda\circ N_{E/\F_{343}}$, where
$\F_{343}=\F_7[t]/(t^3+t^2+1)$ and $\lambda(t+2)=e^{2\pi i/18}$; the
element $t+2$ is primitive.

\subsection{The array}
Index rows and columns by four border indices $\infty_0,\dots,\infty_3$ and
the pairs $(b,x)$ with $b\in\{0,1,2,3\}$ and $x\in E$; the set
$\{b\}\times E$ is the $b$-th \emph{field block}. Let
\[
L=\begin{pmatrix}1&1&1&1\\-1&1&-1&1\\-1&1&1&-1\\-1&-1&1&1\end{pmatrix},\qquad
U=\begin{pmatrix}-1&-1&-1&-1\\-1&-1&1&1\\-1&1&1&-1\\-1&1&-1&1\end{pmatrix}.
\]
Given four patterns $m_0,m_1,m_2,m_3$, define
\begin{equation}\label{eq:array}
\begin{aligned}
H(\infty_a,\infty_c)&=L_{ac}, & H(\infty_a,(c,y))&=U_{ac},\\
H((b,x),\infty_c)&=-U_{cb}, & H((b,x),(c,y))&=L_{bc}\,m_{b\oplus c}\bigl(\phi_{bc}(x,y)\bigr),
\end{aligned}
\end{equation}
where $\oplus$ is bitwise exclusive or, $\phi_{bb}(x,y)=y-x$,
$\phi_{bc}(x,y)=-(x+y)$ if $b\ne c$ and $0\in\{b,c\}$, and
$\phi_{bc}(x,y)=x+y$ otherwise. With $m_0,m_1,m_2,m_3=A,B,C,D$ developed as
$A(x,y)=f_A(y-x)$ and $R$ the permutation matrix of $x\mapsto-x$, the core
of~\eqref{eq:array} is the Wallis--Whiteman array
\[
\begin{pmatrix}
A&BR&CR&DR\\
-BR&A&-D^{\mathsf T}R&C^{\mathsf T}R\\
-CR&D^{\mathsf T}R&A&-B^{\mathsf T}R\\
-DR&-C^{\mathsf T}R&B^{\mathsf T}R&A
\end{pmatrix},
\]
and~\eqref{eq:array} is equation~(11) of~\cite{Kline18}.

An \emph{arrangement} is a choice of which pattern of a system occupies each
position $0,1,2,3$, together with a set of positions whose patterns are
\emph{reversed}, $m(z)\mapsto m(-z)$. Reversal fixes the value at $0$ and
negates the other values. The \emph{standard} arrangement puts $A,B,C,D$ in
positions $0,1,2,3$ with no reversal; for system I in characteristic $7$
this is the matrix of~\cite{Kline18}, with the development direction and the
position of the Paley pattern fixed as in its released code. We write $\bar A$ for $A$ reversed.

\begin{lemma}\label{lem:hadamard}
If the autocorrelations $\sum_{y\in E}m_i(y)m_i(y+s)$ of the four patterns
sum to $-4$ at every $s\ne0$, then every arrangement of them gives a skew
Hadamard matrix of order $4(Q+1)$.
\end{lemma}

\begin{proof}
Each pattern is odd off the origin with value $1$ at $0$, so each
development $M$ satisfies $M+M^{\mathsf T}=2I$ and has row sums $1$, and
reversal preserves autocorrelations. The block computation
of~\cite[\S5]{Kline18} uses only these facts and the sum of the four
autocorrelations, not which pattern sits in which position.
\end{proof}

\subsection{The two systems}

\begin{theorem}\label{thm:gauss}
In the setting of \S2.1, for each system the autocorrelations of $f_{w_A}$,
$f_{w_B}$ and $f_{w_C}$ sum to $-3$ at every nonzero shift. Each of the
three takes, at nonzero shifts, the values $-1$, $-1+t$ and $-1-t$ for one
integer $t\neq0$, and attains each of them.
\end{theorem}

The Paley pattern has autocorrelation $-1$ at every nonzero shift, so
Lemma~\ref{lem:hadamard} applies to both systems.

\begin{proof}[Proof sketch]
Let $G_j$ be the Gauss sum of $\chi^j$ with the canonical additive character.
Frobenius invariance gives $G_j=G_{pj}$, and $p\bmod18\in\{7,13\}$ generates
the subgroup $\{1,7,13\}$, so $G_1G_5=G_9^2=-Q$. The norm exponent $(Q-1)/(p-1)$ is $\equiv3d\pmod{18}$, so
$\chi^6$ is trivial and $\chi^3$ is quadratic on $\F_p^\times$, and
$(3/p)=-1$ because $p\equiv7\pmod{12}$. The Hasse--Davenport product
formula~\cite[\S11.3]{BEW} for the cubic and quadratic characters gives
$G_3=G_1^2/G_5$, $G_{15}=G_5^2/G_1$ and, with $\tau=\chi(4)$, a relation
$\bar\delta=-\tau\delta$ for $\delta=(G_1^2-G_5G_9)/G_2$. The quadratic
formula for $\chi^3$ gives $G_3/G_9=G_6/G_{12}$, because $\chi^3(4)=1$. If
$\delta=0$, then $G_3=G_9$, so $G_6=G_{12}$ and $G_6^2=Q$;
but $G_6\in\Q(\zeta_3)$ because $\chi^6$ is trivial on $\F_p^\times$, and
$\sqrt Q$ is irrational. So $\delta\ne0$.

Expanding each pattern in the characters $\chi^u$ ($u$ odd) writes its
autocorrelation at a shift of class $c$ as
$-\tfrac1{81}\sum_{h\text{ even}}\zeta^{hc}\sum_{u\text{ odd}}
C(u)C(h-u)J(\chi^u,\chi^{h-u})$, with $C(u)=\sum_{j=0}^8w_j\zeta^{-uj}$
and $\zeta=e^{2\pi i/18}$; the term $h=0$ equals $-1$. For $h\ne0$ we
substitute $J(\chi^u,\chi^{h-u})=G_uG_{h-u}/G_h$ and group the terms by
monomials in $G_1,G_5,G_9$. An exact computation in $\Z[\zeta]$ shows that for the
sum of the three patterns every group cancels, and that for each single
pattern every nonconstant term is an explicit integer multiple of $\delta$,
for each of the three values of $\tau$. The specialized coefficients at the
eighteen shifts are permutations of $(0,V,-V)$ with $V\ne0$, and each of
the three patterns attains $0$, $V$ and $-V$; see \S\ref{sec:verify}.
\end{proof}

\begin{remark}
For system I and $p=7$, Theorem~\ref{thm:gauss}, including the three
values, is proved in~\cite{Kline18} by a different route; system II and the
primes $p\ne7$ are new here. The proof above holds
for every primitive character of order $18$.
\end{remark}

\subsection{\texorpdfstring{Lifting from $\F_{343}$}{Lifting from F343}}
The certificates for characteristic $7$ use cyclotomic numbers. For
$Q=343^d$ let $T_d(i,j)$ count $z\in E\setminus\{0,1\}$ with $j(z)=i$ and
$j(1-z)=j$.

\begin{lemma}\label{lem:lift}
Let $d$ be odd. Then $T_d$ equals the $d$-fold cyclic convolution of $T_1$
on $(\Z/18\Z)^2$ plus the constant table with entry $(343^d-2-341^d)/324$.
Consequently, for a prime $\ell\equiv1\pmod{18}$, $T_d\bmod\ell$ depends only
on $d\bmod(\ell-1)$.
\end{lemma}

\begin{proof}
At $(u,v)\ne(0,0)$ the Fourier coefficient of $T_d$ is the Jacobi sum
\[
 J_d(u,v)=\sum_{z\ne0,1}\chi^u(z)\chi^v(1-z).
\] By the Davenport--Hasse
lifting theorem for Gauss sums~\cite[Theorem 5.14]{LidlNiederreiter}
and the expression of a Jacobi sum as $G(\chi^u)G(\chi^v)/G(\chi^{u+v})$,
we have $J_d(u,v)=(-1)^{d-1}J_1(u,v)^d=J_1(u,v)^d$ when $\chi^u$, $\chi^v$
and $\chi^{u+v}$ are nontrivial; in the degenerate cases $J_d$ equals $-1$ or
$-\chi^u(-1)$ for every $d$, and $\chi(-1)=-1$ for every $d$, so again
$J_d=J_1^d$. At $(0,0)$ the coefficients are $Q-2$ and, for the convolution
power, $341^d$. The two tables therefore have the same Fourier transform,
which proves the identity. For the last sentence, reduce modulo a prime
above $\ell$, which maps $\Z[\zeta]$ onto $\F_\ell$; since $324=18^2$ is
invertible modulo $\ell$, $T_d\bmod\ell$ is the inverse transform of the
reductions of $J_1(u,v)^d$ at $(u,v)\ne(0,0)$ and of $343^d-2$ at
$(0,0)$. For $d\ge1$ and
$a\in\F_\ell$, $a^d$ depends only on $d\bmod(\ell-1)$.
\end{proof}

We use $\ell=37$ and $\ell=19$. Each certificate below also checks its values
at $d=1$ by direct summation of matrix entries over $\F_{343}$.

\section{Four-row sets}\label{sec:quartets}

For a set $S$ of four rows of $H$, let $\pi_S$ be the entrywise product of
the four rows, and call the columns where $\pi_S$ takes its less frequent
sign the \emph{minority columns} of $S$. The number of minority columns, and
the set of minority columns, are preserved by equivalence: row signs
multiply $\pi_S$ by a constant, and column signs and permutations act on
columns. Call $S$ \emph{exceptional} if it has at most four minority
columns, and \emph{closed} if it has none. Write $\mathcal B$ for the four
border rows and $\mathcal Z=\{(b,0)\}$ for the four origin rows.

\begin{theorem}\label{thm:quartets}
For both systems, every admissible $p$ and $d$, every character of order
$18$ and every arrangement, the exceptional row sets of $H$ are exactly
$\mathcal B$ and $\mathcal Z$, and each has exactly four minority columns.
The minority columns of $\mathcal B$ are the border columns. Those of
$\mathcal Z$ are the four origin columns if an even number of patterns is
reversed, and the border columns otherwise.
\end{theorem}

\begin{proof}
\emph{The two sets.} $\pi_{\mathcal B}$ is read from $L$ and $U$: it is
constant on each field column block, and equals the minority sign exactly on
the four border columns. For $\mathcal Z$, the row $(b,0)$ restricted to
field column block $c$ is $\pm m_{b\oplus c}(\pm y)$; in each system the four
patterns multiply to $1$ off the origin, and a direct evaluation shows that
$\pi_{\mathcal Z}$ is $+1$ on the border columns, $-1$ on the origin columns,
and $(-1)^{\#\text{reversed}}$ on all other columns.

\emph{Centers.} Let $S$ be any other set of four rows. A row $(b,x)$
restricted to field column block $c$ equals
$\pm m_{b\oplus c}(\pm(y-a))$ with \emph{center} $a=x$ if $b=c$ and
$a=-x$ otherwise. Because patterns are odd off the origin, for $y\ne a$ this
is $\pm m_{b\oplus c}(y-a)$. Two rows of $S$ with the same center in a block
come from different field blocks, since otherwise they would coincide; so the
patterns multiplied at one center are distinct. A border row is constant on
each field block. For both systems the products of one, two or three
distinct patterns are nonconstant on $E^\times$ (indeed they have mean of
absolute value at most $1/3$), and the product of all four is $1$. Hence, if
on every field column block every center carries a constant product, then
either $S$ has no field rows, so $S=\mathcal B$, or $S$ consists of one row
from each field block with a common center in every column block. Writing the
rows as $(b,x_b)$, block $0$ gives $x_0=-x_1=-x_2=-x_3$ and block $1$ gives
$-x_0=x_1=-x_2=-x_3$, so $x_0=-x_0$ and all $x_b=0$: $S=\mathcal Z$.

\emph{Large fields.} Otherwise some column block $c$ has $s\le4$ distinct
centers $a_1,\dots,a_s$ with products $g_1,\dots,g_s$, at least one
nonconstant. Expand $g_k$ on $E^\times$ in the characters $\chi^u$ with
normalized coefficients $c_{k,u}$. For $y$ away from the centers,
\[
 \prod_k g_k(y-a_k)=\sum_{u\in(\Z/18\Z)^s}\Bigl(\prod_k c_{k,u_k}\Bigr)
 \chi\Bigl(\prod_k(y-a_k)^{u_k}\Bigr).
\]
For $u\ne0$, the polynomial $\prod_k(y-a_k)^{u_k}$ is not an eighteenth power,
so Weil's bound~\cite[Theorem 5.41]{LidlNiederreiter} gives
$\abs{\sum_{y\in E}\chi(\cdot)}\le(s-1)\sqrt Q$; removing the centers costs
at most $s$ more, and restoring the true values at the centers at most $s$.
A product of $k$ patterns has the parity $(-1)^k$ under $j\mapsto j+9$, so at
most nine of its coefficients are nonzero and $\sum_u\abs{c_{k,u}}\le3$ by
Parseval and the Cauchy--Schwarz inequality; a single pattern of either system has at most seven nonzero
coefficients, so $\sum_u\abs{c_{k,u}}\le\sqrt7$. If $s=4$, each center carries
a single pattern, whose mean is $0$, and the block sum of $\pi_S$ has absolute
value at most $49(3\sqrt Q+4)+4$. If $s\le3$, the nonconstant center has mean
of absolute value at most $1/3$ and the bound is $Q/3+27(2\sqrt Q+3)+3$. Both
are less than $Q-8$ for $Q\ge31^3$. But four or fewer minority columns force
an absolute block sum of at least $Q-8$ on every field block. The only
admissible field with $Q<31^3$ is $Q=343$.

\emph{The field $\F_{343}$.} Here an exhaustive computation over all row
quartets settles the claim, for each system, for the standard arrangement and
the arrangement with positions $2$ and $3$ exchanged, each with and without
reversal of the pattern in position $0$, and for both orientations
$\chi,\chi^{-1}$ of the character. Every primitive character of order $18$ is
$\chi\circ\varphi$ or $\chi^{-1}\circ\varphi$ for a field automorphism
$\varphi$, and the matrix built from $\chi\circ\varphi$ is obtained from the
one built from $\chi$ by applying $\varphi$ to every field coordinate,
because $\varphi$ is additive and $\phi_{bc}$ is linear; so the two matrices
are equivalent. The computation
uses that two disjoint pairs of rows $\{r_1,r_2\},\{r_3,r_4\}$ form an
exceptional quartet only if the products $r_1r_2$ and $r_3r_4$ agree up to
sign outside four columns, hence on one of five disjoint samples of $32$
columns; all colliding pairs on each sample are then compared in full. The
other arrangements are equivalent to these by Proposition~\ref{prop:transport}
below, through equivalences that map $\mathcal B$ to $\mathcal B$ and
$\mathcal Z$ to $\mathcal Z$. Only the existence of these equivalences is used
here: they are explicit, and checking them does not depend on this
theorem.
\end{proof}

\begin{corollary}\label{cor:closed}
In the matrices of Theorem~\ref{thm:quartets} there are no closed row
quartets. There are exactly two closed column quartets if an even number of
patterns is reversed, and exactly one otherwise. Hence no such matrix is
equivalent to its transpose.
\end{corollary}

\begin{proof}
For a skew matrix and a set $T$ of four indices, the product of the four
columns in $T$ equals the product of the four rows in $T$ off $T$ and its
negative on $T$: off $T$ transposition negates all four factors, on $T$ only
three. So a closed column quartet is a row quartet whose minority columns are
its own four indices. By Theorem~\ref{thm:quartets} this holds for
$\mathcal B$ always and for $\mathcal Z$ exactly when an even number of
patterns is reversed. Exceptional row quartets have four minority columns,
so none is closed.
\end{proof}

\begin{corollary}\label{cor:kron}
No matrix of Theorem~\ref{thm:quartets} is equivalent to a Kronecker product
$H_1\otimes H_2$ of Hadamard matrices of orders greater than $1$, or to a
skew doubling $\bigl(\begin{smallmatrix}K&K\\-K^{\mathsf T}&K^{\mathsf T}
\end{smallmatrix}\bigr)$ of a skew Hadamard matrix $K$.
\end{corollary}

\begin{proof}
In $H_1\otimes H_2$ the rows indexed $(i,k),(i,l),(j,k),(j,l)$ have constant
product, a closed row quartet. In a skew doubling of order $N$, for indices
$i\ne j$ the rows $i$, $j$ of the top half and of the bottom half have product
$-1$ exactly at the coordinates $i$ and $j$ of each half, because
$K^{\mathsf T}$ reverses every off-diagonal sign of a skew $K$; this gives
$\binom{N/2}2>2$ exceptional quartets.
\end{proof}

\section{From an equivalence to one semilinear map}\label{sec:frame}

Fix $p$, $d$ and $\chi$, and let $H$ and $H'$ be matrices of the array for
arrangements of systems I or II. An equivalence $H\to H'$ maps exceptional
quartets to exceptional quartets. We call it \emph{marked} if it maps the
border rows of $H$ to the border rows of $H'$; then it also maps $\mathcal Z$
to $\mathcal Z'$.

\begin{proposition}\label{prop:frame}
A marked equivalence $H\to H'$ maps each field row block and each field
column block of $H$ onto one of $H'$, the four Paley blocks onto the four
Paley blocks, and acts on field coordinates by
\[
(b,x)\mapsto(\pi(b),\varepsilon_bM(x)),\qquad(c,y)\mapsto(\rho(c),\varepsilon'_cM(y)),
\]
with permutations $\pi,\rho$, signs $\varepsilon_b,\varepsilon'_c$, and one
map $M(x)=\beta x^{p^k}$, $\beta\ne0$. Its row and column signs are constant
on each field block.
\end{proposition}

\begin{proof}
\emph{Blocks.} The border columns are the minority columns of $\mathcal B$,
so they are mapped to border columns. On them, a field row $(b,x)$ has signs
$(-U_{cb})_c$, which depend only on $b$, and the four columns of $U$ are
pairwise non-proportional; so the projective sign pattern on the border
columns recovers the field row blocks. Dually the border rows recover the
field column blocks. Comparing an entry in a border column with its image
shows that row signs are constant on each field row block, and dually for
columns.

\emph{Paley blocks.} Inside a field block $(b,c)$, two rows $(b,x),(b,x')$
have inner product equal to the autocorrelation of $m_{b\oplus c}$ at
$x-x'$; reversal does not change it. For the Paley pattern this is $-1$ at
every nonzero shift. By Theorem~\ref{thm:gauss} each other pattern takes both
values $-1\pm t$ with $t\neq0$, and one of them has absolute value different
from $1$. So the Paley blocks are the field blocks in which all row pairs have
inner product $\pm1$, and the equivalence permutes them.

\emph{Semiaffine maps.} In a Paley block the entry is a fixed sign times
$\eta(\pm(y-u))$, extended by $1$ at $0$, where $\eta$ is the quadratic
character and $u=\pm x$. In each row, $(Q-1)/2$ entries have one sign and
$(Q+1)/2$ the other; call the first the minority entries. Since $-1$ is a
nonsquare, after replacing $x$ and $y$ by $\pm x,\pm y$ the minority entries
are exactly the incidences $y-x\in S$ of the Paley design, with $S$ the
nonzero squares. Row and column signs are constant on the block, so the
equivalence maps the block's entries to the target block's entries up to one
sign, and the counts force minority entries to minority entries. Thus the
equivalence induces permutations $F,G$ of $E$ with $F(y)-G(x)\in S$ whenever
$y-x\in S$.
Kantor's Corollary~8.2 in~\cite{Kantor2T}, valid for odd prime powers
greater than $11$ (here $Q\equiv3\pmod4$, the case of his Theorem~8.1), gives $F=G$ and $F(x)=\beta x^\sigma+\alpha$ with
$\beta\in S$ and $\sigma\in\Aut(E)$. Each field row block and each field
column block lies in exactly one Paley block, so all eight coordinate maps
are of this form up to the substitutions $x\mapsto\pm x$. The origin row of
each block is mapped to an origin row, so the row maps fix $0$; since $F=G$,
so do the column maps. Hence every coordinate map is additive.

\emph{One map.} In a field block write the argument of the pattern as
$h(y+tx)$ with $h,t\in\{\pm1\}$ ($t=-1$ on the diagonal). Let $P,S'$ be the
additive maps on the row and column coordinates of this block, and let $t'$
belong to the target block. Holding $z=y+tx$ fixed, the image argument is
$h'(S'(z)+(t'P-tS')(x))$, and the target entry does not depend on $x$.
Hence every element of the image of the additive map $t'P-tS'$ is a period of
the target pattern. A nonzero additive subgroup $W$ of $E$ has order a power
of $p$, and a function constant on its cosets has sum divisible by $p$; but
every pattern sums to $1$. So $t'P=tS'$ for every block. Since every row
block meets every column block, all eight maps are $\pm M$ for one $M$.
\end{proof}

A marked equivalence is therefore described by finite data: the
permutations of the border rows, the field row blocks, the border columns and
the field column blocks, the signs, and the action of $M$ on classes,
\[
j(\pm M(x))=h+p^kj(x),\qquad h\in\Z/18\Z,\ p^k\bmod18\in\{1,7,13\}.
\]
Call such data a \emph{template}. The template determines whether entries
agree on every class and at every origin, and it depends on $p$ only through
$p\bmod18\in\{7,13\}$, which generate the same group $\{1,7,13\}$.
Conversely, a template that preserves all entries is realized by each of the
$(Q-1)/18$ elements $\beta$ of the class $h$ (after the sign) and each of the
$d$ exponents $k'\in\Z/3d\Z$ with $p^{k'}\equiv p^k\pmod{18}$, giving
exactly $\Lambda_Q$ distinct marked equivalences per template.

\begin{proposition}\label{prop:transport}
For each system, the following holds for every admissible $p$, $d$ and
$\chi$. Two arrangements are related by a marked equivalence if and only if
their permutations of the four patterns have the same sign and their numbers
of reversed patterns have the same parity. Each of the four resulting classes
contains $96$ of the $384$ arrangements. The marked self-equivalences of an
arrangement form $1$ template for system I and $3$ templates for system II.
\end{proposition}

\begin{proof}
A finite enumeration over all templates. With the sign of the first border
row normalized, the program runs through the $24^4$ permutations of border
rows, field row blocks, border columns and field column blocks; the
$8\times8$ submatrix on border and origin indices then determines all row and
column signs, $96$ choices preserve it, and $48$ of these are compatible with
the relations $t'P=tS'$. For each of the $48$ it tries the $54$ class actions
$(h,p^k\bmod18)$ and compares the patterns on all eighteen classes and at the
origin in each of the sixteen field blocks. For system II the three
self-templates have $h=0,3,6$ (and $0,12,15$ after reversal of the pattern
in position $0$), all with $p^k\equiv1\pmod{18}$.
\end{proof}

\section{Automorphism groups}\label{sec:aut}

\begin{theorem}\label{thm:groups}
Let $H$ be any arrangement of system I or II for admissible $p$, $d$, $\chi$,
and let $\Aut_0(H)$ be the group of row-and-column permutations underlying its
automorphisms, i.e.\ automorphisms modulo negation of all rows and columns.
\begin{enumerate}
\item The subgroup of $\Aut_0(H)$ fixing the set $\mathcal B$ has order $\Lambda_Q$
for system I and $3\Lambda_Q$ for system II, and $\Aut_0(H)$ has order $\Lambda_Q$ or
$2\Lambda_Q$ for system I, and $3\Lambda_Q$ or $6\Lambda_Q$ for system II.
\item For the standard arrangement of system I, the subgroup fixing
$\mathcal B$ is the group of maps $x\mapsto ax^{p^{3j}}$, $a\in(E^\times)^{18}$,
applied to every field copy, fixing all border indices; it is
$K\rtimes\Gal(E/\F_{p^3})$ with $K=(E^\times)^{18}$.
\item Let $p=7$. Every automorphism fixes $\mathcal B$ for every
arrangement of system I when $\chi$ is the character of~\cite{Kline18}, and
for every arrangement of system II with an odd number of reversed patterns
and every character $\chi$ of order $18$.
\end{enumerate}
\end{theorem}

Theorem~\ref{thm:aut} is part (2) with part (3); the order is $19$ at
$Q=343$, and the group is cyclic there because $K$ is cyclic and $d=1$.

\begin{proof}
(1) Proposition~\ref{prop:transport} and the count of $\Lambda_Q$ realizations per
template give the stabilizer of $\mathcal B$. An automorphism permutes the
two exceptional quartets, so this stabilizer has index at most $2$.

(2) The single self-template of the standard arrangement of system I has all
permutations trivial, all signs positive, $h=0$ and $p^k\equiv1$.

(3) Mark either exceptional quartet, delete it and its four minority columns,
and partition the remaining rows by their projective signs on the deleted
columns and the remaining columns by their projective signs on the deleted
rows. This gives sixteen submatrices, and the number of them in which every
two rows have inner product $\pm1$ is an invariant of the marking. For the
border marking the submatrices are the field blocks, and exactly four, the
Paley blocks, pass. For the origin marking, a certificate shows that none of
the sixteen passes: for each it exhibits a border row and a field row of the
same part whose restricted inner product is not $\pm1$ modulo $37$. The
inner products are linear in the cyclotomic numbers $T_d$, so
Lemma~\ref{lem:lift} with $\ell=37$ reduces all odd $d$ to the eighteen
residues of $d$ modulo $36$, and each value is also checked at $d=1$ by
summing matrix entries. The certificate covers the four arrangements of
system I with positions $2,3$ standard or exchanged and position $0$
reversed or not, for the character of~\cite{Kline18}, and the two
arrangements of system II with position $0$ reversed, for both orientations
$\chi,\chi^{-1}$; every character of order $18$ is the image of one of these
under a field automorphism, which gives an equivalent matrix (see the proof of Theorem~\ref{thm:quartets}). By
Proposition~\ref{prop:transport} these represent all arrangements in
question.
\end{proof}

\begin{remark}
At $Q=343$ the order $19$ was also confirmed by an exhaustive
individualization-and-refinement search for isomorphisms between the
normalized bipartite graphs of $H$, anchored at a border row and border
column, with every accepted map checked against all entries. The search
restricts the image of the anchors to the $32$ pairs allowed by
Theorem~\ref{thm:quartets}, and finds exactly the $19$ scalar maps.
\end{remark}

\section{Equivalence classes}\label{sec:classes}

\begin{proof}[Proof of Theorem~\ref{thm:systems}]
The matrices are skew Hadamard by Lemma~\ref{lem:hadamard} and
Theorem~\ref{thm:gauss}. By Corollary~\ref{cor:closed}, the ordered pair
(closed row quartets, closed column quartets) is $(0,2)$ for the standard
arrangements, $(0,1)$ after reversing $A$, and $(2,0)$ and $(1,0)$ for their
transposes. By Theorem~\ref{thm:groups}(1), the automorphism group order
lies in $\{\Lambda_Q,2\Lambda_Q\}$ for system I and in $\{3\Lambda_Q,6\Lambda_Q\}$ for system II,
and transposition preserves it. These two invariants separate all eight
matrices.
\end{proof}

\begin{proof}[Proof of Theorem~\ref{thm:sixteen}]
By Proposition~\ref{prop:transport} each system has at most four classes of
arrangements and four of transposes, and the transposes are distinct from the
arrangements by Corollary~\ref{cor:closed}. Systems I and II are separated
by automorphism group order. For system I, every equivalence between
arrangements is marked by Theorem~\ref{thm:groups}(3), so the four classes
of Proposition~\ref{prop:transport} are distinct: eight classes with the
transposes. For system II, the closed column quartets separate even from odd
reversal parity, and the two odd classes are distinct because every
equivalence between them is marked. It remains to compare the two even
classes, represented by the arrangements $(A,B,C,D)$ and $(A,B,D,C)$.

For these two matrices, delete the eight rows of $\mathcal B\cup\mathcal Z$
and their eight minority columns. Group the remaining columns by their
projective signs on the eight deleted rows; there are $24$ cells, each the
union of three classes of one field block. For a choice $T\in\{\mathcal B,
\mathcal Z\}$, group the $24$ cells into four groups of six by their
projective signs on $T$. For each remaining row $r$ and each group, form the
six absolute cell sums $\abs{\sum_{c\in C}H_{rc}H_{sc}}$ with an anchor $s$
among the deleted rows (any anchor gives the same value), reduce them modulo
$19$ up to sign, and sort them. Let $n_T(\tau)$ be the number of pairs (row,
group) with sorted tuple $\tau$. The unordered pair
$\{n_{\mathcal B}(\tau),n_{\mathcal Z}(\tau)\}$ is invariant under all
equivalences, including those that exchange the two quartets. By
Lemma~\ref{lem:lift} with $\ell=19$ these counts depend only on $d\bmod18$;
they are multiples of $(Q-1)/18$. Table~\ref{tab:hist} gives, for each
residue with $3\nmid d$, a tuple whose pair of counts differs between the two
matrices. For $d\equiv3,9,15\pmod{18}$ the pairs agree for every tuple, and
the test is inconclusive. Hence system II gives four classes of arrangements
when $3\nmid d$ and three or four when $3\mid d$, and twice as many with
transposes.
\end{proof}

\begin{table}[ht]
\centering
\begin{tabular}{cccc}
\toprule
$d\bmod18$ & tuple $\tau$ & $(A,B,C,D)$ & $(A,B,D,C)$\\
\midrule
1  & $(1,5,5,5,7,7)$ & $\{72,72\}$ & $\{0,72\}$\\
5  & $(1,1,1,3,6,8)$ & $\{72,144\}$ & $\{72,72\}$\\
7  & $(1,3,3,4,4,7)$ & $\{0,0\}$ & $\{0,72\}$\\
11 & $(1,3,3,4,4,7)$ & $\{0,72\}$ & $\{0,0\}$\\
13 & $(1,1,1,3,6,8)$ & $\{72,72\}$ & $\{72,144\}$\\
17 & $(1,5,5,5,7,7)$ & $\{0,72\}$ & $\{72,72\}$\\
\bottomrule
\end{tabular}
\caption{Witness tuples for the two even arrangements of system II. Counts
are divided by $(Q-1)/18$.}\label{tab:hist}
\end{table}

\begin{proof}[Proof of Theorem~\ref{thm:compare}]
Transposes, Kronecker products and skew doublings are excluded by
Corollaries~\ref{cor:closed} and~\ref{cor:kron}. Suppose $H$ is equivalent
to a Hadamard matrix $H^\ast$ of the array whose four functions are constant
on the cosets of the sixth powers. Multiplication of every field coordinate
by a generator of the sixth powers is an automorphism of $H^\ast$ of odd order
$\mu=(Q-1)/6$; it fixes the eight border and origin rows and columns, and all
its other cycles have length $\mu$. Its conjugate $g$ is an automorphism of $H$
with the same cycle type. Being of odd order, $g$ fixes each exceptional
quartet and its minority columns as sets, and these sets have fewer than
$\mu$ elements, so $g$ fixes them pointwise. Then $g$ lies in the stabilizer of
$\mathcal B$ and fixes every border row and border column. Since the columns
of $U$ are pairwise non-proportional, the border rows then show that $g$
maps each field column block to itself, and dually each field row block. So
its template has trivial permutations. Then the border and origin entries
force all row and column signs to agree, and the relations $t'P=tS'$ force
all coordinate signs to agree; with the normalization of
Proposition~\ref{prop:transport}, the class condition reads $m(h+p^kj)=m(j)$
for each of the four patterns $m$ of the system and every class $j$, and a
reversed pattern gives the same condition. That condition does not depend on
the arrangement, and the enumeration of Proposition~\ref{prop:transport} for
the standard arrangement shows that only $h=0$, $p^k\equiv1\pmod{18}$
satisfies it. So $g$ has the form
$x\mapsto ax^{p^{3j}}$ with $a\in(E^\times)^{18}$. If the field automorphism is trivial, the order of $g$
divides $(Q-1)/18$, which $\mu$ does not. If it has order $b\ge3$, then $g^b$
is multiplication by an element of a subfield with $p^{3d/b}$ elements, so the
order of $g$ is at most $b(p^{3d/b}-1)<\mu$. Either way the order is not
$\mu$.
\end{proof}

\section{Verification}\label{sec:verify}

Every finite step above is an exact computation; Table~\ref{tab:certs}
lists them with the programs, which are in the directory \texttt{code} of the
repository. They need only Node.js and a C++17 compiler; together they run in
about two and a half minutes on one core, and the repository pins every
output.

\begin{table}[ht]
\centering\small
\begin{tabular}{p{0.42\textwidth}p{0.5\textwidth}}
\toprule
Step & Check\\
\midrule
Theorem~\ref{thm:gauss}: grouped Gauss-sum identities and profiles; Fourier
supports and subset means used in Theorem~\ref{thm:quartets} &
\texttt{families/uniform.cjs}: exact arithmetic in $\Z[\zeta_{18}]$, both
systems, all three values of $\tau$; cross-checked in explicit fields of sizes
$7^3,31^3,43^3$ by \texttt{base/family-crosscheck.cjs}\\
Lemma~\ref{lem:hadamard} at $Q=343$ & \texttt{base/gram.cpp}:
$HH^{\mathsf T}=NI$ and $H+H^{\mathsf T}=2I$ for all $768$ arrangements\\
Theorem~\ref{thm:quartets} at $Q=343$ & \texttt{algebra/flip-quartets.cjs}:
all exceptional row quartets, both systems, the representative arrangements,
both orientations of $\chi$\\
Proposition~\ref{prop:transport} & \texttt{algebra/skeleton.cjs},
\texttt{families/second-skeleton.cjs}, \texttt{algebra/placements.cjs}\\
Theorem~\ref{thm:groups}(3) & \texttt{algebra/frame.cjs},
\texttt{families/odd-frames.cjs}: origin marking modulo $37$, eighteen
residues of $d\bmod36$, direct checks at $d=1$\\
Table~\ref{tab:hist} & \texttt{families/even-frames.cjs}: histogram modulo
$19$, nine odd residues of $d\bmod18$, direct checks at $d=1$\\
Remark after Theorem~\ref{thm:groups} & \texttt{base/aut.cpp}: exhaustive
automorphism search at order $1376$\\
\bottomrule
\end{tabular}
\caption{Exact certificates.}\label{tab:certs}
\end{table}

\section{Remarks}\label{sec:remarks}

\emph{Scope.} The theorems compare the matrices of this construction with
each other and with the families named in Theorem~\ref{thm:compare}. They do
not compare them with every other skew Hadamard matrix of the same orders,
such as those of Goethals--Seidel type, and they give no new orders.

\emph{Open cases.} When $3\mid d$ we do not know whether the two even
arrangements of system II are equivalent; an equivalence would have to
exchange the border rows with the origin rows. Outside the cases of
Theorem~\ref{thm:groups}(3), that is, for $p\ne7$, for system I with the
other orientation of the character, and for system II with an even number of
reversed patterns, we do not know whether any automorphism exchanges the two
exceptional quartets, which is why Theorem~\ref{thm:groups}(1) gives two
possible orders. The sixteen-class
count uses one character; the other orientation of the character was checked
for the quartet theorem and for the origin-marking certificate of system II,
but not for that of system I or for the histogram.

\subsection*{Role of AI}
AI assisted the derivations, the computer verification, the literature
search and the exposition under the author's direction. Process-separated AI
reviews checked the proofs and certificates; they are not independent expert
peer review, and no independent human expert has examined the proofs. The
release repository records the sources, the reviews and their dispositions.

\begin{thebibliography}{9}
\bibitem{BEW}
B.~C. Berndt, R.~J. Evans and K.~S. Williams,
\emph{Gauss and Jacobi Sums}, Wiley, New York, 1998.

\bibitem{Hall62}
M.~Hall, Jr.,
\emph{Note on the Mathieu group $M_{12}$},
Arch. Math. \textbf{13} (1962), 334--340.

\bibitem{Kantor69}
W.~M. Kantor,
\emph{Automorphism groups of Hadamard matrices},
J. Combin. Theory \textbf{6} (1969), 279--281.
\href{https://doi.org/10.1016/S0021-9800(69)80088-8}{doi:10.1016/S0021-9800(69)80088-8}.

\bibitem{Kantor2T}
W.~M. Kantor,
\emph{2-transitive symmetric designs},
Trans. Amer. Math. Soc. \textbf{146} (1969), 1--28.
\href{https://doi.org/10.1090/S0002-9947-1969-0255422-6}{doi:10.1090/S0002-9947-1969-0255422-6}.

\bibitem{Karoui}
A.~Karoui,
\emph{An explicit skew-Hadamard matrix of order 1252 via cyclotomic unions},
arXiv:2602.16089, 2026.

\bibitem{Kline18}
J.~Kline,
\emph{Order-18 Jacobi triples beyond the sextic shift orbit}, version 0.1.0,
2026.
\href{https://doi.org/10.5281/zenodo.21783219}{doi:10.5281/zenodo.21783219}.

\bibitem{Lehmer74}
E.~Lehmer,
\emph{A family of supplementary difference sets},
Bull. Austral. Math. Soc. \textbf{11} (1974), 1--4.
\href{https://doi.org/10.1017/S0004972700043550}{doi:10.1017/S0004972700043550}.

\bibitem{LidlNiederreiter}
R.~Lidl and H.~Niederreiter,
\emph{Finite Fields}, second edition,
Cambridge University Press, 1997.

\bibitem{MomiharaXiang}
K.~Momihara and Q.~Xiang,
\emph{Skew Hadamard difference families and skew Hadamard matrices},
arXiv:1801.08776, 2018.

\bibitem{Orrick}
W.~P. Orrick,
\emph{Switching operations for Hadamard matrices},
SIAM J. Discrete Math. \textbf{22} (2008), 31--50.
\href{https://doi.org/10.1137/050641727}{doi:10.1137/050641727}.

\bibitem{Pantangi}
V.~R.~T. Pantangi,
\emph{Difference families, skew Hadamard matrices, and critical groups of
doubly regular tournaments}, Electron. J. Combin. \textbf{26} (2019), no.~3,
Paper P3.59.
\href{https://doi.org/10.37236/8753}{doi:10.37236/8753}.

\bibitem{WallisWhiteman}
J.~S. Wallis and A.~L. Whiteman,
\emph{Some classes of Hadamard matrices with constant diagonal},
Bull. Austral. Math. Soc. \textbf{7} (1972), 233--249.
\href{https://doi.org/10.1017/S0004972700044993}{doi:10.1017/S0004972700044993}.

\bibitem{Yamada92}
M.~Yamada,
\emph{Supplementary difference sets and Jacobi sums},
Discrete Math. \textbf{103} (1992), 75--90.
\href{https://doi.org/10.1016/0012-365X(92)90041-D}{doi:10.1016/0012-365X(92)90041-D}.
\end{thebibliography}
\end{document}
